\scp{\tsc{\ili{Teor}-Kur}}{11-04}
The \ili{Teor} language is spoken on an island to the south of the
\ili{Watubela} string of islands,
and \ili{Kur} is spoken on the nearby northwestern most islands of the \ili{Kei} archipelago.
Data is extremely scarce for both languages making their precise
relationship to other languages difficult to determine.
Nonetheless, what data is available suggests they belong within STB.
This classification must be revisited when more data for these languages becomes available.
Both languages were tentatively linked (with reservations) to
\ili{Kei} and Tanimbar languages by \citet{Hughes1987} based on lexical
similarity that is closer to \ili{Kei} (54{\%})
and \ili{Fordata} (57{\%}) than to \tsc{Eastern Islands} languages of Maluku.
But lexical similarity with \ili{Yamdena} (48{\%}) is about the same
as with \ili{Geser} (49{\%}).

Regarding the data sources,
the only accessible data for \ili{Teor} is the word list in \citet{Wallace1869}.\footnote{
		Shortly before the final version of this manuscript was
		submitted for typesetting and publication we received a 479-item
		\ili{Teor} wordlist (see appendix \ref{ch:ComWor}) from Craig Marshall.
		This wordlist confirms that \ili{Teor} has the defining sound changes
		of \tsc{Seram-Tanimbar-Bomberai} with *d/*z > \it{r} and
		*ə/*a/*-ay > \it{a} in final syllables (there is insufficient for *-aw).
		We have not had time to update this section according to this new data.}
The only data for \ili{Kur} we have been able to access is the word list
in \citet[108--114]{vonRosenberg1867},\footnote{
		The \ili{Kur} word list
		in \citet{vonRosenberg1867} is reproduced with a
		different orthography in \citet[600--613]{vonRosenberg1878}.
		\citet{vonRosenberg1867,vonRosenberg1867} also has a word list
		marked “Tijoor” but the language of this word list is a variety
		of \ili{Geser}-\ili{Gorom}.}
as well as two folktales,
\citet{Ohoiwutun2021} and \citet{Rado2021},
published by the Indonesian language office.
\ili{Teor} data from \citet{Wallace1869} is given as transcribed in
that source within angle brackets ‹ ›.
Similarly, \ili{Kur} data from \citet{vonRosenberg1867} is given in angle brackets.
The transcription follows Dutch orthographic practices,
for which the following phonetic interpretations are secure:
‹j› = [y], ‹oe› = [u], and ‹ie› = [i].
\ili{Kur} data from \citet{Ohoiwutun2021}
and \citet{Rado2021} follows Indonesian orthography.
We present this data without any brackets,
only re-transcribing ‹ng› as ‹ŋ›.

\ili{Teor} and \ili{Kur} appear to share *z/*d > **d with STB,
though data is extremely scarce.
All reflexes of *z
and *d we can identify are shown in \qf{11-04}.
Reflexes of *quzan ‘rain’ have *z > \it{r},
thus merging with reflexes of *d.
However, reflexes of *zaqat have a different reflex with *z > \it{y}.
Some members of the \tsc{Tanimbar-Bomberai} node also have *z
> \it{y} in this etymon (see \trf{11-03})
and this may suggest some kind of link
between \tsc{\ili{Teor}-Kur} and \tsc{Tanimbar-Bomberai}.
However, more reflexes of *z are clearly needed before any definitive
statements can be made.
\ili{Teor} *daRaq > ‹larah› blood may also attest a link between
\tsc{\ili{Teor}-Kur} and \tsc{Tanimbar-Bomberai} as members of \tsc{Nuclear Tanimbar-Bomberai}
(\srf{11-05-01}) also have *d > (**r >) \it{l} /{\gap}R through dissimilation.

\noindent{\wg\st{6pt}\tblr{>{\hspace{-\tabcolsep}}l<{\hspace{-\tabcolsep}}
>{\hspace{-\tabcolsep}\enspace}
llll}
\lsptoprule
%\tbx{11-4}	&	PMP	&	\ili{Teor}	&	\ili{Kur}	&	local gl.	\\	\midrule
%\endfirsthead
\tbx{11-04}	&	PMP	&	\ili{Teor}	&	\ili{Kur}	&	local gl.	\\	\midrule
%\endhead
	&	*\tbr{d}uha	&	‹\tbr{r}ua›	&	\tbr{r}ua	&	two	\\
	&	*si-i\tbr{d}a	&		&	hi\tbr{r}a	&	they	\\
	&	*ma-u\tbr{d}əhi	&		&	mui\tbr{r}	&	go behind	\\
	&	*\tbr{d}aRaq	&	‹\tbr{l}arah›	&		&	blood	\\
	&	*qu\tbr{z}an	&	‹hu\tbr{r}ani›	&	‹oe\tbr{r}an›	&	rain	\\
	&	*\tbr{z}aqat	&	‹\tbr{y}at›	&	‹\tbr{j}aaid›	&	bad	\\
\lspbottomrule
\end{tabular}\mbox{}\\}

The other defining sound change of STB
is *ə/*a/*-ay/*-aw > \it{a} in final syllables.
Every identifiable reflex of *a,
*ə, and *-ay/*-aw in final syllables for \ili{Teor}
and \ili{Kur} is shown in \qf{11-05}--\qf{11-08}.
The reflexes of final *-aw in *qaləjaw ‘day,
sun’ in \qf{11-07} may be irregular,
perhaps due to the influence of adjacent *j
before *j > {\0}, or they may show that
\ili{Teor} and \ili{Kur} do not have *-aw > \it{a}.

\noindent{\wg\st{6pt}\tblr{>{\hspace{-\tabcolsep}}l<{\hspace{-\tabcolsep}}
>{\hspace{-\tabcolsep}\enspace}
lllll}
\lsptoprule
%\tbx{11-05}	&	PMP	&	\ili{Teor}	&	\ili{Kur}	&	local gl.\\	\midrule
%\endfirsthead
&	&	PMP	&	\ili{Teor}	&	\ili{Kur}	&	local gl.\\	\midrule
%\endhead
\tbx{11-05}	&*ə&	*ma-qit\tbr{ə}m	&	‹mit\tbr{e}n›	&	‹mid mid\tbr{a}n›	&	black	\\
	&&	*tan\tbr{ə}m	&		&	da-tan\tbr{a}n	&	plant (v.)	\\
	&&	*tan\tbr{ə}q	&		&	kantan\tbr{a}, ‹tan-tan\tbr{a}›	&	earth, land	\\
	&&	*dəŋ\tbr{ə}R	&		&	da-nan\tbr{a}r	&	hear	\\
	&&	*nip\tbr{ə}n	&	‹nif\tbr{i}n›	&	nip\tbr{a}n	&	tooth	\\
	&&	*ən\tbr{ə}m	&	‹nem›	&	‹nèèn›	&	six	\\	\midrule
\end{tabular}\mbox{}\\}

\noindent{\wg\st{6pt}\tblr{>{\hspace{-\tabcolsep}}l<{\hspace{-\tabcolsep}}
>{\hspace{-\tabcolsep}\enspace}
lllll}
\lsptoprule
%\tbx{11-05}	&	PMP	&	\ili{Teor}	&	\ili{Kur}	&	local gl.\\	\midrule
%\endfirsthead
&	&	PMP	&	\ili{Teor}	&	\ili{Kur}	&	local gl.\\	\midrule
%\endhead
\tbx{11-06}	&*-ay&	*maRuqan\tbr{ay}	&	‹menránn\tbr{a}›	&	manran\tbr{a}	&	man	\\
	&&	*mat\tbr{ay}	&		&	i-mat\tbr{a}	&	die	\\
	&&	*paj\tbr{ay}	&	‹pas\tbr{er}›	&	pas\tbr{a}	&	rice	\\
	&&	*qaq\tbr{ay}	&	‹ya\tbr{i}n›	&	‹ja\tbr{i}n›	&	leg, foot	\\
\lspbottomrule
\end{tabular}\mbox{}\\}
	
\noindent{\wg\st{6pt}\tblr{>{\hspace{-\tabcolsep}}l<{\hspace{-\tabcolsep}}
>{\hspace{-\tabcolsep}\enspace}
lllll}
\lsptoprule
%\tbx{11-05}	&	PMP	&	\ili{Teor}	&	\ili{Kur}	&	local gl.\\	\midrule
%\endfirsthead
&	&	PMP	&	\ili{Teor}	&	\ili{Kur}	&	local gl.\\	\midrule
%\endhead
\tbx{11-07}	&*-aw&	*qaləj\tbr{aw}	&	‹lil\tbr{éw}›	&	l\tbr{o}ŋ,‹l\tbr{è}›	&	day	\\
	&&	*qaləj\tbr{aw}	&	‹l\tbr{ew}›	&		&	sun	\\	\midrule
\tbx{11-08}	&*a&	*duh\tbr{a}	&	‹rú\tbr{a}›	&	ru\tbr{a}	&	two	\\
	&&	*daR\tbr{a}q	&	‹lar\tbr{a}h›	&		&	blood	\\
	&&	*quz\tbr{a}n	&	‹hur\tbr{a}ni›	&	‹oer\tbr{a}n›	&	rain	\\
	&&	*bul\tbr{a}n	&	‹phul\tbr{a}n›	&	fu\tbr{a}n	&	moon	\\
	&&	*am\tbr{a}	&	‹am\tbr{a}›	&		&	father	\\
	&&	*əp\tbr{a}t	&	‹f\tbr{a}ht›	&	‹pa\tbr{ä}t›	&	four	\\
	&&	*ma-əs\tbr{a}	&		&	mes\tbr{a}ŋ	&	alone	\\
	&&	*tuq\tbr{a}h	&		&	fitu\tbr{a}	&	old woman	\\
	&&	*Rum\tbr{a}q	&	\cg{(sarin)}	&	rum\tbr{a}	&	house	\\
	&&	*si-id\tbr{a}	&		&	hir\tbr{a}	&	\tsc{3pl}	\\
	&&	*hik\tbr{a}n	&	‹ik\tbr{a}n›	&	ik\tbr{a}n	&	fish	\\
	&&	*in\tbr{a}	&	‹in\tbr{a}›	&	‹in\tbr{ă}n›	&	mother	\\
	&&	*lim\tbr{a}	&	‹lim\tbr{a}›	&	‹lim\tbr{a}›	&	five	\\
	&&	*qasiR\tbr{a}	&	‹sir\tbr{e}n›	&		&	salt	\\
	&&	*siw\tbr{a}	&	‹siw\tbr{e}›	&	‹siw\tbr{a}h›	&	nine	\\
	&&	*miñ\tbr{a}k	&	‹min\tbr{e}k›	&		&	sweet	\\
	&&	*bu\tbr{a}q	&	‹phu\tbr{i}n›	&	fu\tbr{e}n	&	fruit	\\
	&&	*bu\tbr{a}q	&		&	fu\tbr{a}	&	betel nut	\\
	&&	*an\tbr{a}k	&		&	an\tbr{a}k, an\tbr{i}k	&	child	\\
	&&	*wak\tbr{a}R	&	‹wok\tbr{i}›	&	wak\tbr{i}r	&	root	\\
	&&	*qalim\tbr{a}	&	‹lim\tbr{i}n›	&		&	hand	\\
	&&	*taŋ\tbr{a}n	&	‹kimin tag\tbr{i}n›	&		&	finger	\\
	&&	*babaq	&		&	faf\tbr{i}n	&	below	\\
	&&	*mat\tbr{a}	&	‹mat\tbr{i}n›	&	‹mat\tbr{i}n›	&	eye	\\
	&&	*zaq\tbr{a}t	&	‹yat›	&	‹jaaid›	&	bad	\\
\lspbottomrule
\end{tabular}\mbox{}\\}

\ili{As} can be seen from this data,
both languages appear to have *a/*ə/*-ay > \it{i} in final syllables
for words which have a third person possessor.
Some of these words may in fact attest *a/*ə > {\0} /{\gap}{\#}
with subsequent attachment of a suffix \it{-in}.
Whatever the exact process,
\ili{Kur} has variant forms for *anak (\it{an\tbr{ak} hira} ‘children’
and \it{an\tbr{i}k} ‘his/her child’) showing that
some kind of productive morphology is probably involved.
Similarly, words with final \it{u} also have \it{i}
when possessed by a third person,
though these examples may attest *u > \it{i} /{\gap}C{\#}.
One example which can be gleaned from the \ili{Kur} data is *bat\tbr{u}
+ *pəñu > \it{fa\tbr{u}t pen} ‘turtle rock’ but *bat\tbr{u} > \it{fat\tbr{i}n} ‘seed'.
Potential examples from \ili{Teor} are given in \qf{11-00}.

\noindent{\wg\st{6pt}\tblr{>{\hspace{-\tabcolsep}}l<{\hspace{-\tabcolsep}}
>{\hspace{-\tabcolsep}\enspace}
lll}
\lsptoprule
\tbx{11-00}	&	PMP	&	\ili{Teor}	& local gl.	\\	\midrule
	&	*qul\tbr{u}	&	‹ul\tbr{i}n›	&	head	\\
	&	*kuhk\tbr{u}h	&	‹limin kuk\tbr{i}n›	&	fingernail	\\
	&	*mas\tbr{u}	&	‹yaf me\tbr{i}n›	&	smoke	\\
	&	*bul\tbr{u}	&	‹phul\tbr{i}n›	&	hair	\\
	&	*kamb\tbr{u}	&	‹kab\tbr{i}n›	&	belly	\\
\lspbottomrule
\end{tabular}\mbox{}\\}

Apart from these examples,
\ili{Kur} has *ə/*a/*-ay > \it{a}.
\ili{Teor} has *ə/*a/*-ay > ‹a› in most cases,
but *ə/*a > ‹a›, ‹e› /iC{\gap}{\#}.
Thus, \ili{Teor} *ma-qit\tbr{ə}m > ‹mit\tbr{e}n› ‘black’ is consistent
with *siw\tbr{a} > ‹siw\tbr{e}› ‘nine’,
*qasiR\tbr{a} > ‹sir\tbr{e}n› ‘salt’,
and *miñ\tbr{a}k > ‹min\tbr{e}k› ‘sweet’
and we can tentatively posit that *ə
and *a merged in final syllables in \ili{Teor}.

\ili{Kur} is an important witness in the region as it is the only language
which retains *p unchanged in all positions.
\ili{Teor}, like other languages of the region,
has lenition of *p > ‹f›.
Examples of *p in the available data are given in \qf{11-09}.

\noindent{\wg\st{6pt}\tblr{>{\hspace{-\tabcolsep}}l<{\hspace{-\tabcolsep}}
>{\hspace{-\tabcolsep}\enspace}
llll}
\lsptoprule
%\tbx{11-9}	&	PMP	&	\ili{Teor}	&	\ili{Kur}	&	local gl.	\\	\midrule
%\endfirsthead
\tbx{11-09}	&	PMP	&	\ili{Teor}	&	\ili{Kur}	&	local gl.	\\	\midrule
%\endhead
	&	*\tbr{p}əñu	&		&	\tbr{p}en	&	turtle	\\
	&	*\tbr{p}itu	&	‹\tbr{f}it›	&	‹\tbr{p}iet›	&	seven	\\
	&	*\tbr{p}anid	&	‹\tbr{f}anik›	&		&	wing	\\
	&	*ə\tbr{p}at	&	‹\tbr{f}aht›	&	‹\tbr{p}aät›	&	four	\\
	&	*ha\tbr{p}uy	&	‹ya\tbr{f}›	&	yau\tbr{p}, ‹iaa\tbr{p}›	&	fire	\\
	&	*sa-\tbr{p}uluq	&		&	‹sa\tbr{p}oel›	&	ten	\\
	&	*ni\tbr{p}ən	&	‹ni\tbr{f}in›	&	‹ni\tbr{p}an›	&	tooth	\\
\lspbottomrule
\end{tabular}\mbox{}\\}

\ili{Teor} and \ili{Kur} have *j > {\0} in all reflexes which have been identified,
with the exception of reflexes of *pajay ‘rice’.\footnote{
		\ili{Teor} data provided by Craig Marshall has medial *j > \it{l}
		in all reflexes except those of *qaləjaw.}
The retention of *p = \it{p} in \ili{Teor} indicates that this term
is a loan from another language.
Examples of *j in the available data are given in \qf{11-10}.
%*j > {\0} may be a defining change of \tsc{\ili{Teor}-Kur} setting
%it apart from other branches of STB,
%though more data is needed to be confident.

\noindent{\wg\st{6pt}\tblr{>{\hspace{-\tabcolsep}}l<{\hspace{-\tabcolsep}}
>{\hspace{-\tabcolsep}\enspace}
llll}
\lsptoprule
%\tbx{11-10}	&	PMP	&	\ili{Teor}	&	\ili{Kur}	&	local gl.	\\	\midrule
%\endfirsthead
\tbx{11-10}	&	PMP	&	\ili{Teor}	&	\ili{Kur}	&	local gl.	\\	\midrule
%\endhead
	&	*qalə\tbr{j}aw	&	‹liléw›	&	loŋ, ‹lè›	&	day	\\
	&	*qalə\tbr{j}aw	&	‹lew›	&		&	sun	\\
	&	*ŋa\tbr{j}an	&		&	ŋaa-m	&	your name (‹ngam› in \citealp[21]{Rado2021})	\\
	&	*kuni\tbr{j}	&	‹kúni›	&		&	yellow	\\
	&	*pa\tbr{j}ay	&	‹pa\tbr{s}er›	&	pa\tbr{s}a	&	rice	\\
\lspbottomrule
\end{tabular}\mbox{}\\}

PMP *s undergoes a split in \ili{Teor} and \ili{Kur}.
\ili{Teor} has *s = \it{s} in two reflexes,
with one other potential reflex showing *s > {\0}.
\ili{Kur} has *s > \it{s, h}.
Examples are given in \qf{11-11}.\footnote{
		Another potential
		reflex of *s in \ili{Kur} is *asu > ‹yaa› \citep[33]{Rado2021},
		‹eja› \citep{vonRosenberg1867}.
		Given Indonesian orthographical conventions in which a sequence
		of two identical vowels has a medial glottal stop,
		‹yaa› is probably [yaʔa] which can be compared with \ili{Kei}
		\it{yaha} and \ili{Fordata} \it{yahaw} ‘dog’ reflecting intermediate **yasaw,
		which may be from *asu with irregular final *u > \it{aw}
		and expected \it{y-}insertion before vowel-initial words,
		see \trf{11-28}.} 

\noindent{\wg\st{5pt}\tblr{>{\hspace{-\tabcolsep}}l<{\hspace{-\tabcolsep}}
>{\hspace{-\tabcolsep}\enspace}
llll}
\lsptoprule
%\tbx{11-11}	&	PMP	&	\ili{Teor}	&	\ili{Kur}	&	local gl.	&		\\	\midrule
%\endfirsthead
\tbx{11-11}	&	PMP	&	\ili{Teor}	&	\ili{Kur}	&	local gl.		\\	\midrule
%\endhead
	&	*\tbr{s}iwa	&	‹\tbr{s}iwer›	&	‹\tbr{s}iwah›	&	nine			\\
	&	*\tbr{s}a-puluq	&		&	‹\tbr{s}apoel›	&	ten			\\
	&	*qa\tbr{s}iRa	&	‹\tbr{s}iren›	&		&	salt			\\
	&	*ma-ə\tbr{s}a	&		&	me\tbr{s}aŋ	&	alone			\\
	&	*ə\tbr{s}a	&		&	\tbr{h}a	&	one			\\
	&	*\tbr{s}i-ida	&		&	\tbr{h}ira	&	\tsc{3pl}			\\
	&	*qa\tbr{s}awa	&		&	\tbr{h}oin	&	his wife			\\
	&	**ma-qasu	&	‹yaf mein›	&		&	smoke	(*s > {\0},
																			*u > \it{i} /σ\und{σ}{\#};
																			‹yaf› = fire)	\\

\lspbottomrule
\end{tabular}\mbox{}\\}

\largerpage
PMP *b is reflected as ‹f›,
‹ph›, or ‹w› in \ili{Teor}.
Orthographic ‹ph› may represent a bilabial fricative [ɸ]
as opposed to labio-dental [f] ‹f›.
\ili{Kur} data from \citet{Ohoiwutun2021}
and \citet{Rado2021} has *b > f.
Whether this ‹f› represents a labio-dental fricative or
bilabial fricative is unknown.
Data from \citet{vonRosenberg1867} has *b > ‹w› in three
words and *b > ‹f› in one.
Whether the difference in the data from \citet{vonRosenberg1867}
and that in \citet{Ohoiwutun2021}
and \citet{Rado2021} is a dialect difference, a change which
has occurred in the last 150 years, or only a transcription difference is unknown.
Examples of PMP *b are given in \qf{11-12}.

\noindent{\wg\st{5.5pt}\tblr{>{\hspace{-\tabcolsep}}l<{\hspace{-\tabcolsep}}
>{\hspace{-\tabcolsep}\enspace}
lllll}
\lsptoprule
%\tbx{11-12}	&	PMP	&	\ili{Teor}	&	\ili{Kur}	&	\ili{Kur} \citep{vonRosenberg1867}	&	local gl.	\\	\midrule
%\endfirsthead
\tbx{11-12}	&	PMP	&	\ili{Teor}	&	\ili{Kur}	&	\ili{Kur}	&	local gl.	\\ \hhline{~}
&&&&	\citep{vonRosenberg1867}	&	\\	\midrule
%\endhead
	&	*\tbr{b}a\tbr{b}aq	&		&	\tbr{f}a\tbr{f}in	&		&	below	\\
	&	*\tbr{b}atu	&		&	\tbr{f}atin	&		&	seed, granule	\\
	&	*\tbr{b}atu	&		&	\tbr{f}aut	&		&	stone	\\
	&	*\tbr{b}anua	&		&	\tbr{f}inua	&	‹\tbr{w}anoea›	&	village	\\
	&	*\tbr{b}uaq	&		&	\tbr{f}ua	&		&	betel nut	\\
	&	*\tbr{b}uaq	&	‹\tbr{ph}uin›	&	\tbr{f}uen	&		&	fruit	\\
	&	*\tbr{b}ulan	&	‹\tbr{ph}ulan›	&	\tbr{f}uan	&	‹\tbr{w}oea› 	&	moon	\\
	&	*\tbr{b}uRiq	&		&	i-\tbr{f}urik	&		&	wash (rice)	\\
	&	*\tbr{b}a\tbr{b}uy	&	‹\tbr{f}a\tbr{f}›	&		&	‹\tbr{f}a\tbr{f}›	&	pig	\\
	&	*\tbr{b}ulu	&	‹\tbr{ph}ulin›	&		&		&	feather	\\
	&	*\tbr{b}ahi-tuqah	&		&	\tbr{f}itua	&		&	grandmother	\\
	&	*ma-\tbr{b}\<in\>ahi	&	‹me\tbr{w}ina›	&	ma\tbr{f}ina	&	‹men\tbr{w}iena›	&	woman	\\
	&	*\tbr{b}\<in\>ahi	&	‹we\tbr{w}ina›	&	\cg{(hoin)}	&		&	wife	\\
\lspbottomrule
\end{tabular}\mbox{}\\}

PMP *k > \it{k, h} in \ili{Teor}, while 
\ili{Kur} has *k = \it{k}.
Examples are given in \qf{11-13}.

\noindent{\wg\st{6pt}\tblr{>{\hspace{-\tabcolsep}}l<{\hspace{-\tabcolsep}}
>{\hspace{-\tabcolsep}\enspace}
llll}
\lsptoprule
%\tbx{11-13}	&	PMP	&	\ili{Teor}	&	\ili{Kur}	&	local gl.	\\	\midrule
%\endfirsthead
\tbx{11-13}	&	PMP	&	\ili{Teor}	&	\ili{Kur}	&	local gl.	\\	\midrule
%\endhead
	&	*\tbr{k}ulit	&	‹\tbr{h}olit›	&		&	skin	\\
	&	*\tbr{k}utu	&	‹\tbr{h}ut›	&		&	louse	\\
	&	*hi\tbr{k}an	&	‹i\tbr{k}an›	&	i\tbr{k}an	&	fish	\\
	&	*\tbr{k}ambu	&	‹\tbr{k}abin›	&		&	belly	\\
	&	*\tbr{k}ayu	&	‹\tbr{k}ai›	&	\tbr{k}ai	&	wood	\\
	&	*\tbr{k}ahu	&		&	\tbr{k}a	&	\tsc{2sg}	\\
	&	*\tbr{k}aən	&		&	i-\tbr{k}an	&	eat	\\
	&	*\tbr{k}amuyu	&		&	\tbr{k}im	&	\tsc{2pl}	\\
	&	*\tbr{k}u\tbr{k}u	&	‹limin \tbr{k}u\tbr{k}in›	&		&	fingernail	\\
	&	*\tbr{k}a\tbr{k}a	&		&	\tbr{k}a\tbr{k}a-muir	&	siblings	\\
	&	{\#}mu\tbr{k}u	&	‹mū\tbr{k}›	&	‹mu\tbr{k}›	&	banana	\\
	&	*wa\tbr{k}aR	&	‹wo\tbr{k}i›	&	wa\tbr{k}ir	&	root	\\
	&	*ta\tbr{k}ut	&		&	ammata\tbr{k}ut	&	scared	\\
	&	*a\tbr{k}u	&		&	au\tbr{k}	&	1sg	\\
	&	*manu\tbr{k}	&	‹mano\tbr{k}›	&	manu\tbr{k}	&	bird	\\
	&	*miña\tbr{k}	&	‹mine\tbr{k}›	&		&	sweet	\\
	&	*ana\tbr{k}	&		&	ana\tbr{k}, ani\tbr{k}	&	child	\\
\lspbottomrule
\end{tabular}\mbox{}\\}

PMP *ŋ > \it{g} in \ili{Teor},
though there is only one example; *ta\tbr{ŋ}an > ‹limin ta\tbr{g}in› ‘finger’.
The change is confirmed by \ili{Kur} \it{li\tbr{ŋ}ain}, \ili{Teor} ‹la\tbr{g}ain›
‘path, way’, as well as \citet[84]{Hughes1987} who states; “Also,
velar nasals in \ili{Kur} become voiced velar plosives in Ut [\ili{Teor}].”
The change *ŋ > \it{g} is regionally common and also occurs
in nearby \ili{Kowiai} (\srf{11-03-02} and \tsc{East Bomberai} (\srf{11-05-02}).
Denasalisation (*ŋ > \it{k} /V{\gap}V) also occurs in \tsc{East Seram} (\srf{11-02}).

\citet[84]{Hughes1987} also observes one difference between \ili{Kur}
and his \ili{Teor} data (collected from Ut village in the \ili{Kei} Islands):
“There are some consistent phonetic changes between \ili{Kur}
and Ut, particularly the palatalisation in Ut of word-final alveolar
nasals, laterals and voiceless alveolar plosives in \ili{Kur}.” Such palatalisation
is not observed in the \ili{Teor} data in \citet{Wallace1869}
and may be due to not being observed/transcribed by the collector
of the list, a dialect difference, or a change
which occurred since \citeauthor{Wallace1869}’s
list was collected.

%no. (sort)	core	no.	Indonesian	note	target language	target note	PMP	
%5	*	5	tali hutan				*waRəj	woras
%75	*	73	nama				*ŋajan	gali-n
%77	*	75	adik (kelamin yg sama)				*huaji	ak walug
%102	*	102	matahari				*qaləjaw	leuʷ
%103	*	103	hari				*qaləjaw	hadin le
%111	*	112	hujan				*quzan	uran
%133	*	134	gunung				*bukij, **letay	fukar
%137	*	138	arang				*qajəŋ	yaran
%178	*	175	lalat				*laləj, *bəRŋaw	kalbu:r
%187	*	184	ulat				*quləj	(kapapao)
%252		247	dekat				*zaŋi, *daŋi	(her)
%254	*	249	jauh				*zauq	loinʸ
%264	*	259	pahit				*(ma-)paqit, **ma-pəju	(geran)
%267	*	262	jahat				*zaqat	yač
%268	*	263	(jalan) kering				*maja	(sʸer)
%301	*	295	dagu				*qazay	(ka gekum)
%307	*	301	hidung				*ijuŋ, *ŋijuŋ	gilinʸ
%320	*	314	tunjuk				*tuzuq	ia pyaturuk
%322	*	316	jengkal				*zaŋkal	pagat
%335	*	329	pusar				*pusəj	pulinʸ
%338	*	332	empedu				*qapəju	pelinʸ
%342	*	336	sakit				*hapəjəs, *hapəjiq	(migat)
%348	*	342	ludah				**kambeR, *qizuR	ia kaberinʸ
%356	*	350	jalan (setapak)				*zalan	laganʸ
%364	*	358	tangga				*haRəzan	leran
%382	*	376	tajam				*tazəm	masneran
%389	*	383	jarum				*zaRum	lanlen
%390	*	384	jahit				*zaqit, **sa(q)u	ia hyor